\chapter{Committing Reality}

\section{Two Different Kinds of "Done"}

Chapter 14 ended with Y, the nearest point in C to a proposed update, computed and ready. It is tempting to treat this as the end of the story: a proposal was made, it was checked, it was corrected where necessary, and the result is simply what the field becomes next. This chapter argues that the temptation should be resisted, because computing Y and the field actually becoming Y are two different kinds of event, and collapsing them erases a distinction this book's roadmap insisted on from the very first sketch of its core cycle back in Chapter 6.

Y is a fact about mathematics. Given \(M_t\), \(\Delta\psi\), and C, the projection \(\Pi_C\)(\(M_t\) + \(\Delta\psi\)) has a determinate answer, computable in principle without anything actually having happened to the world. \(M_{t+1}\) is a fact about history. It is what the field genuinely becomes, the state that Chapter 11's residue will remember having occurred, the state future observations will be checked against. A Y that is computed but never committed is not a diminished version of \(M_{t+1}\). It is not part of the field's history at all, in the same sense that a sentence drafted and then deleted was never part of the document it might have joined. Commit is the operation that turns the first kind of done into the second, and it deserves to be treated as its own event rather than an automatic consequence of projection having finished.

\section{Why Projection Alone Cannot Decide}

There is a concrete reason projection cannot simply be trusted to settle the matter on its own, beyond the conceptual tidiness of keeping computation and history separate. Projection, as Chapter 14 developed it, answers a question about a single proposal in isolation: given this \(\Delta\psi\), what is the nearest admissible state? It does not, on its own, know whether some other proposal, arising elsewhere or from another agent at the same moment, has already claimed the very state Y would move toward, or proposed something incompatible with it. Two proposals can each individually project to a perfectly admissible Y, and still be jointly impossible to commit together, because their combination is not itself admissible even though neither alone was at fault. Something has to arbitrate this, and it cannot be projection, because projection was never given the other proposal to check against in the first place.

\section{Multiple Proposals, One History}

This becomes unavoidable once more than one agent is proposing changes to the same field at overlapping moments, which any persistent, shared world must eventually allow. Two kinds of relationship can hold between simultaneous proposals. Updates are commutative when they touch disjoint or otherwise compatible parts of the field, one agent's proposal at the courtyard's bench, another's proposal at a doorway on the far side, in which case both may commit without either needing to know about the other or without their order mattering to the final result. Updates are non-commutative when they conflict, both proposing incompatible changes to the same region, in which case order or arbitration matters and at most one, or some admissible combination of parts of each, can actually be committed.

This book will not resolve conflicts of the second kind through an explicit negotiation protocol, a designated authority deciding whose proposal wins. Resolution instead follows admissibility directly: among the proposals contending for a given region, whichever combination remains inside C after joint projection is what commits, and proposals or portions of proposals that cannot be reconciled with the rest are the ones left out. Committed changes accumulate into a causal record, a structure this book will develop no further here than to note that it behaves as a directed graph of dependency rather than a single linear sequence, since two commutative updates genuinely have no fact of the matter about which happened first. Chapter 29 returns to this in its fully engineered form, as the concrete mechanism behind sharing one persistent world across multiple simultaneous observers and agents; this chapter's task is only to establish that commit, not projection, is where multiplicity gets resolved.

\section{Refusal Revisited}

Chapter 14 closed by flagging a case its own machinery could not fully handle: a proposal so incompatible with the field's accumulated state that the nearest admissible point in C is not really a correction of it at all, but something else entirely, honoring almost none of what was actually proposed. That chapter left the question open because projection, as a pure nearest-point computation, has no mechanism for declining to answer. It always returns some Y, however distant.

Commit does have such a mechanism, and this is where this book locates it. Before a computed Y is written into the field's history, commit may decline to write it, treating the original proposal as void rather than silently substituting the distant admissible point Y happens to name. This is refusal in the sense this book's supporting material gestures at without fully developing: not a correction performed on a proposal, but an irreversible decision that a given branch of possibility, the one the proposal was reaching for, does not become part of this field's history at all. A proposal that gets refused leaves no trace, no partial commitment, no compromise state standing in for what was actually intended. It simply does not happen, and the field proceeds as though it had never been proposed. This book will not specify, at this stage, exactly how large a distance between proposal and projection triggers refusal rather than acceptance; that threshold depends on details Part IV is better positioned to formalize. What matters here is that refusal has a home now, at commit rather than at projection, and that the question Chapter 14 left open was never really about projection's mathematics. It was about a decision projection was never equipped to make.

\section{Commit as the Point of No Return}

This gives residue, introduced in Chapter 11, its precise relationship to the write cycle this Part has been developing. Residue accumulates only from committed events. A Y that was computed, considered, and refused contributes nothing to \(R_{t+1}\), exactly as though the proposal behind it had never been generated. A Y that is committed becomes part of the field's history permanently, in the sense this book has used permanence throughout: not immune to future change, a wall can still be damaged, a bench can still be moved, but no longer erasable as though it had never occurred. This is what makes commit the point of no return in this framework, and it is the reason this book's roadmap insisted, several chapters before this one was drafted, that projection and commit be kept as visibly separate steps rather than folded into a single operation. A projected state is a possibility under consideration. A committed state is a fact the rest of the field's future is now obligated to be consistent with.

\section{A First Glimpse of Stability}

One consequence of this chapter's account is worth naming now, briefly, ahead of the chapter that develops it properly. A field that keeps committing genuinely new content, region by region, cannot remain forever in flux; some regions will eventually stop changing, not because nothing more could be proposed for them, but because whatever continues to be proposed there projects and commits to the same state it already held. A courtyard whose bench has settled, whose residue is dense and specific, and whose entropy has long since fallen to near zero is a region where \(\Delta\psi\), once selected and projected, tends to commit to something indistinguishable from what was already there. This is the first appearance, informal and unstated as a rule, of the fixed-point behavior Chapter 16 exists to formalize.

\section{Looking Ahead}

This chapter has separated computing an admissible state from actually committing it, explained why multiple simultaneous proposals require commit rather than projection to arbitrate between them, and given refusal, left unresolved at the end of Chapter 14, a proper home as a decision commit is entitled to make. What remains is to ask what happens when this cycle, propose, project, commit, repeats indefinitely at a single region without anything further changing. That stability, and the conditions that produce or fail to produce it, is the subject of Chapter 16.
